\section{Metrics and Assessment Criteria}
\label{sec:metrics}

We will use a number of different metrics to assess the performance of
the submitted algorithms. Many of these metrics, as well as the
motivations behind them, are defined and discussed in
Ref.~\cite{therailteam2025}.

\subsection{Metrics for per-object point estimates}

Performance on per-object point estimates, i.e., providing a single
best estimate of the redshift of each object in the test sample. All
of our point-estimate metrics first compute the scaled residual,
$\Delta_i$, between the point estimate for each object, $z_{p,i}$, and
the true redshift for that object, $z_{t,i}$:

\begin{equation}
    \Delta_i = \frac{z_{p,i}-z_{t,i}}{1+z_{t,i}}.
\end{equation}

We then use this to construct the following metrics; see
Ref.~\cite{therailteam2025} for more details:

\begin{itemize}
    \item \texttt{Bias} is simply the median of $\Delta_i$.
    \item \texttt{OutlierRate} is the fraction of the $\Delta$ distribution
      outside of $[0, {\rm max}(0.06, 3\sigma_{\rm iqr})]$.
    \item \texttt{SigmaMAD} is an estimate of the standard deviation of
      the median absolute deviation (MAD), which is computed as      
    \begin{equation}
    {\sigma_{\rm MAD}} = 1.4862\,{\rm median}(|\Delta_i - {\rm median}(\Delta)|).
    \end{equation}
\end{itemize}


\subsection{\texorpdfstring{Metrics for per-object $p(z)$ distributions}{Metrics for per-object p(z) distributions}}

We will also assess the algorithm's ability to provide a precise and
accurate estimate of the posterior distribution, $p(z)$, for each
object using the following metrics.

\begin{itemize}
    \item{Conditional Density Loss (\texttt{CDELoss}): We implement
      the method in \cite{2017arXiv170408095I} to compute this metric.
      A better estimation should return a smaller CDE loss.}
    \item{Probability Integral Transform (PIT): This is the
        cumulative distribution function of the photo-$z$ PDF evaluated
        at the galaxy's true redshift for each galaxy in the catalog,
        i.e., ${\rm PIT} = \int_{0}^{z_t} p(z)\, dz$.
      Following Ref.~\cite{schmidt_evaluation_2020}, we provide the PIT-QQ
      (quantile-quantile) diagram, where the PIT distribution is
      directly compared to the ideal uniform distribution. 
      A diagonal PIT-QQ diagram indicates a good estimation.          
      We will then use three metrics to quantify how well the PIT
      distribution matches the ideal: the Kolmogorov-Smirnov (KS)
      test, the Root Mean Square Error (RMSE), and the
      Kullback-Leibler divergence (KL divergence).}
\end{itemize}


\subsection{Metrics for computational usability and performance}

We will assess relevant aspects of the computational performance that
will affect usability and scaling. 

\begin{itemize}
\item{\textbf{Ease of use}: We will assess whether the algorithm is easy to
    install and can be run on the different task sets without needing
    excessively complicated additional configuration files.}
\item{\textbf{Training time}: How quickly the algorithm trains models,
    and how this scales with the training sample size. Here we
    mainly want to ensure that the training time will not dominate the
    iteration cycle. Taking several minutes to train on 100k objects
    is fine; taking hours to do so would be problematic.}
\item{\textbf{Model size}: How large the trained model files are,
    and how this scales with the training sample size. Again, we
    mainly want to ensure that the model size will not tax our resources.
    If the model files are an order of magnitude larger than the input
    data files, we might worry.}
\item{\textbf{Estimation time}: How quickly the algorithm estimates
    redshifts per object. This will determine the use cases for
    which we might use the algorithm. We can run an algorithm that
    takes a few ms per object on all of the billions of galaxies we
    will have in the final LSST sample; for an algorithm that takes a
    few seconds per object, we would probably be constrained to only
    run it on much smaller particular datasets for specific science
    cases, such as samples of supernovae or strongly lensed objects.}
\item{\textbf{Output data size per object}: How large the output
    files with the $p(z)$ estimates are. For a \texttt{qp}
    interpolated grid representation with 300 points, these
    would be about 2.4 kB per object, which is large but manageable,
    whereas for a Gaussian mixture model with 5 Gaussians, this would
    be close to 120 bytes per object.}
\end{itemize}

%\subsection{Metrics for bin assignment}

%\eac{Efficiency and purity}


%\subsection{Metrics for ensemble $n(z)$ estimation}

%\eac{Bias, width, simplified fisher forecast}


%\subsection{Metrics for ensemble estimation performance}

%\eac{Training time, per-object estimation time, model size}

